\section{Discussion}

The results support a careful version of the circuit hypothesis. Toxic and detoxified
text occupy linearly separable regions of representation space in two multilingual
decoder models across multiple fine-tuning seeds. The transfer gap shows that this
signal is not uniformly language-invariant. Exploratory steering sweeps show that
interventions can modify output behavior in a strength-sensitive way without obvious
mode collapse. They do not, by themselves, prove that a shared causal detoxification
circuit has been isolated.

This distinction matters for multilingual safety research. Preservation metrics quantify
whether a rewrite stays close to an input or reference, but cannot substitute for a
valid non-toxicity measure. Likewise, a probe can reveal a useful representation signal
without showing that the signal is sufficient for controllable generation. The
controlled intervention and human-evaluation stages are necessary for interpreting
steering as detoxification rather than generic output perturbation.
